\documentclass[11pt,a4paper]{article}

\usepackage[utf8]{inputenc}
\usepackage[T1]{fontenc}
\usepackage[english]{babel}
\usepackage[a4paper,margin=2.2cm]{geometry}
\usepackage{amsmath,amssymb,amsfonts}
\usepackage{bm}
\usepackage{mathtools}
\usepackage{enumitem}
\usepackage{booktabs}

\setlist[itemize]{leftmargin=1.6em}
\setlength{\parindent}{0pt}
\setlength{\parskip}{0.55em}
\setcounter{tocdepth}{2}

\begin{document}

\tableofcontents

\section{Notation and overall structure}

At each timepoint $t$, the dataset contains a cell point cloud
\[
\mathcal{C}_t=\{c_1,\dots,c_{N_t}\},
\]
with, for each cell, a 3D position
\[
\mathbf{x}_i(t)=
\begin{pmatrix}
x_i(t)\\
y_i(t)\\
z_i(t)
\end{pmatrix}\in\mathbb{R}^3.
\]

The pipeline computes:
\begin{itemize}
  \item a persistent biological identity \texttt{unique\_cell\_id};
  \item raw coordinates $(x,y,z)$ and stabilized coordinates $(x_{\mathrm{stab}},y_{\mathrm{stab}},z_{\mathrm{stab}})$;
  \item migration metrics;
  \item neighborhood relationships;
  \item an embryonic surface at each timepoint;
  \item a local density defined on the vertices of that surface.
\end{itemize}

\section{Assignment of cell identities and lineage events}

\paragraph{Intuitive principle.}
This step reconstructs a biologically coherent identity from successive detections. The problem is formulated as a matching problem between cells observed at two consecutive timepoints. The core algorithm is the Hungarian algorithm, also known as the Kuhn--Munkres algorithm, applied to a distance matrix. Differences in cardinality between two consecutive timepoints are then interpreted as mitoses, disappearances, or fusions according to simple geometric rules.

Assignment is performed independently within each \texttt{track\_id}. If a file does not contain a \texttt{track\_id}, we set
\[
\texttt{track\_id} \leftarrow \texttt{cell\_id}.
\]
If some rows do not have a \texttt{track\_id}, a distinct synthetic identifier is created for each of them.

For two consecutive timepoints $t_{k-1}$ and $t_k$, we consider
\[
P=\{\mathbf{p}_1,\dots,\mathbf{p}_{n_{\mathrm{prev}}}\},
\qquad
Q=\{\mathbf{q}_1,\dots,\mathbf{q}_{n_{\mathrm{curr}}}\},
\]
where $\mathbf{p}_i$ are the positions of the current cells at the previous timepoint and $\mathbf{q}_j$ are the detections at the current timepoint.

The cost matrix is the Euclidean distance matrix
\[
C_{ij}=\|\mathbf{p}_i-\mathbf{q}_j\|_2.
\]

\subsection{Case $n_{\mathrm{curr}}=n_{\mathrm{prev}}$}

Matching is obtained with the Hungarian algorithm (Kuhn--Munkres):
\[
\pi^\star=\arg\min_{\pi}\sum_{i=1}^{n_{\mathrm{prev}}} C_{i,\pi(i)}.
\]

If $C_{i,\pi^\star(i)}>60$, the association is rejected and a new cell is created; otherwise, the previous identity is propagated.

\subsection{Case $n_{\mathrm{curr}}>n_{\mathrm{prev}}$: mitosis}

An initial Hungarian matching is first computed between the previous cells and a subset of the current detections. Unmatched detections are interpreted as candidate new daughter cells.

For each unmatched detection $\mathbf{q}_{\mathrm{new}}$, the mother cell is chosen among the already matched previous cells by minimizing
\[
d_{\mathrm{mother}}=\min\Big(
\|\mathbf{q}_{\pi(i)}-\mathbf{q}_{\mathrm{new}}\|_2,\,
\|\mathbf{p}_{i}-\mathbf{q}_{\mathrm{new}}\|_2
\Big).
\]

The mother cell is then replaced by two new daughter identities. The mother cell is marked as a mitotic event at the previous timepoint.

\subsection{Case $n_{\mathrm{curr}}<n_{\mathrm{prev}}$: disappearance or fusion}

All current detections are matched to a subset of the previous cells. Unmatched previous cells are then tested for a possible fusion event: if their closest current detection is at distance $<60$, they are fused into the same daughter cell.

If several cells $\{i_1,\dots,i_m\}$ fuse into the same detection $j$, a new identity is created with
\[
\texttt{parent\_cell}=i_1/\cdots/i_m.
\]

Otherwise, a large distance $C_{ij}>60$ is interpreted as the appearance of a new independent cell.

\section{Coordinate stabilization}

\paragraph{Intuitive principle.}
The purpose of stabilization is to remove the global motion of the embryo so that measured displacements reflect cellular motions as accurately as possible. The method assumes that, from one image to the next, a subset of cells acts as a sufficiently stable internal reference frame. Registration between two timepoints is therefore treated as a rigid optimal superposition problem between two sets of homologous points. The algorithm used is the Kabsch algorithm, i.e. the SVD solution of the orthogonal Procrustes problem in 3D.

\subsection{Selection of reference cells}

Stabilization is computed from a subset of reference cells:
\begin{itemize}
  \item present at at least two timepoints;
  \item not involved in mitosis, neither as mother nor as daughter.
\end{itemize}

For each consecutive pair $(t_{k-1},t_k)$, only reference cells present at both timepoints are retained. If their number is less than $3$, no stabilization is applied for that transition.

\subsection{Rigid alignment problem}

Let $\mathcal{R}_k$ be the set of reference cells common to $t_{k-1}$ and $t_k$. We define
\[
P=
\begin{bmatrix}
\mathbf{p}_1^\top\\
\vdots\\
\mathbf{p}_m^\top
\end{bmatrix},
\qquad
Q=
\begin{bmatrix}
\mathbf{q}_1^\top\\
\vdots\\
\mathbf{q}_m^\top
\end{bmatrix},
\]
where:
\begin{itemize}
  \item $\mathbf{p}_i$ is the \emph{already stabilized} position of cell $i$ at $t_{k-1}$;
  \item $\mathbf{q}_i$ is its raw position at $t_k$.
\end{itemize}

We seek the rigid transformation $(R,\mathbf{t})$ minimizing
\[
\min_{R,\mathbf{t}} \sum_{i=1}^{m}\left\|
\mathbf{p}_i-(R\mathbf{q}_i+\mathbf{t})
\right\|_2^2,
\qquad
R\in SO(3).
\]

\subsection{Kabsch solution}

We define the centroids
\[
\bar{\mathbf{p}}=\frac{1}{m}\sum_{i=1}^{m}\mathbf{p}_i,
\qquad
\bar{\mathbf{q}}=\frac{1}{m}\sum_{i=1}^{m}\mathbf{q}_i.
\]

The centered coordinates are
\[
\tilde{\mathbf{p}}_i=\mathbf{p}_i-\bar{\mathbf{p}},
\qquad
\tilde{\mathbf{q}}_i=\mathbf{q}_i-\bar{\mathbf{q}}.
\]

The cross-covariance matrix is
\[
H=\tilde{Q}^\top \tilde{P}.
\]

Its SVD is
\[
H=USV^\top.
\]

The optimal rotation is
\[
R=V
\begin{pmatrix}
1&0&0\\
0&1&0\\
0&0&\operatorname{sign}\big(\det(VU^\top)\big)
\end{pmatrix}
U^\top,
\]
which enforces $\det(R)=+1$ and avoids a reflection.

The translation is
\[
\mathbf{t}=\bar{\mathbf{p}}-R\bar{\mathbf{q}}.
\]

\subsection{Application to all cells}

The transformation estimated on reference cells is applied to \emph{all} cells at timepoint $t_k$:
\[
\mathbf{x}_{\mathrm{stab}}(t_k)=R\,\mathbf{x}(t_k)+\mathbf{t}.
\]

The algorithm is sequential: each timepoint is aligned onto the previous \emph{already stabilized} timepoint. Stabilization therefore results from a composition of successive rigid transformations.

\subsection{Error quantification}

Before alignment:
\[
\mathrm{RMSE}_{\mathrm{before}}
=
\sqrt{
\frac{1}{m}
\sum_{i=1}^{m}
\|\mathbf{p}_i-\mathbf{q}_i\|_2^2
}.
\]

After alignment:
\[
\mathrm{RMSE}_{\mathrm{after}}
=
\sqrt{
\frac{1}{m}
\sum_{i=1}^{m}
\|\mathbf{p}_i-(R\mathbf{q}_i+\mathbf{t})\|_2^2
}.
\]

The maximum residual is
\[
r_{\max}=
\max_{1\le i\le m}
\|\mathbf{p}_i-(R\mathbf{q}_i+\mathbf{t})\|_2.
\]

The script also stores the cumulative rotation as Euler angles $(x,y,z)$ for diagnostic purposes.

\section{Migration metrics}

\paragraph{Intuitive principle.}
Once coordinates are stabilized, each trajectory is summarized by a small set of geometric and kinematic descriptors. These quantities measure the total path length, the net displacement, the deviation from a straight trajectory, and the instantaneous velocities.

For a cell observed at timepoints $t_1<\dots<t_n$, with positions $\mathbf{x}(t_k)$, we define:

\[
L=\sum_{k=1}^{n-1}\|\mathbf{x}(t_{k+1})-\mathbf{x}(t_k)\|_2
\]
as the path length,

\[
D=\|\mathbf{x}(t_n)-\mathbf{x}(t_1)\|_2
\]
as the net displacement,

\[
S=\frac{D}{L}
\]
as the straightness index if $L>0$.

If the time interval between two images is $\Delta t$ minutes, the instantaneous velocity on interval $[t_k,t_{k+1}]$ is
\[
v_k=
\frac{\|\mathbf{x}(t_{k+1})-\mathbf{x}(t_k)\|_2}
{(t_{k+1}-t_k)\Delta t}.
\]

The script reports the mean velocity
\[
\bar{v}=\frac{1}{n-1}\sum_{k=1}^{n-1} v_k
\]
and the maximum velocity
\[
v_{\max}=\max_k v_k.
\]

\section{Cell neighborhood}

\paragraph{Intuitive principle.}
Neighborhood is not defined as a simple independent computation at each image. It is initialized by a $k$ nearest-neighbor search, then propagated over time while explicitly accounting for lineage events. The goal is to track biologically coherent neighborhood relationships rather than a purely instantaneous neighborhood. The algorithmic core is therefore a \emph{k-nearest neighbors} (k-NN) search, here with $k=3$, enriched by inheritance rules for mitosis and fusion.

\subsection{Definition of initial neighbors}

Neighborhood is computed on stabilized coordinates by default. For each cell $c$, we consider its first appearance timepoint $t_0(c)$. Among all cells active at the same timepoint, we select the $N=3$ nearest neighbors in the Euclidean sense:
\[
\mathcal{N}_0(c)=
\operatorname*{arg\,min}_{\substack{\mathcal{S}\subset \mathcal{C}_{t_0(c)}\setminus\{c\}\\|\mathcal{S}|=3}}
\sum_{n\in\mathcal{S}}
\|\mathbf{x}_c(t_0)-\mathbf{x}_n(t_0)\|_2.
\]

In practice, this amounts to sorting the distances
\[
d(c,n;t_0)=\|\mathbf{x}_c(t_0)-\mathbf{x}_n(t_0)\|_2
\]
and retaining the three smallest.

\subsection{Inheritance during lineage events}

Neighborhood is not recomputed from scratch at each image. It is \emph{inherited} and updated.

\paragraph{Mitosis.}
If a daughter cell $c'$ originates from a mother cell $m$, it inherits the neighbors of $m$:
\[
\mathcal{N}_0(c') \leftarrow \mathcal{N}_0(m),
\]
then the following are removed:
\begin{itemize}
  \item $c'$ itself;
  \item the other daughters of the same mother.
\end{itemize}

If this leaves fewer than three neighbors, the list is completed with the closest active cells at the daughter cell's appearance timepoint.

\paragraph{Fusion.}
If a cell $c'$ results from the fusion of parents $m_1,\dots,m_r$, the ordered union of parental neighborhoods is taken:
\[
\mathcal{N}_0(c') \leftarrow
\bigcup_{j=1}^{r} \mathcal{N}_0(m_j),
\]
after removing duplicates, the parents themselves, and $c'$.

\subsection{Temporal evolution of neighborhood}

For each cell $c$, a time-dependent neighborhood list is constructed, denoted $\mathcal{N}(c,t)$.

\paragraph{If a neighbor undergoes mitosis at time $t$.}
If $n\in \mathcal{N}(c,t^-)$ divides into daughters $n_1,n_2$, then
\[
n \;\mapsto\; \{n_1,n_2\}.
\]

\paragraph{If several neighbors fuse at time $t$.}
If a subset of neighbors fuses into a cell $n^\star$, then all these neighbors are replaced by $n^\star$.

\subsection{Reported distances}

For each cell $c$, each timepoint $t$ at which it exists, and each active neighbor $n\in\mathcal{N}(c,t)$, the script computes
\[
d(c,n;t)=\|\mathbf{x}_c(t)-\mathbf{x}_n(t)\|_2.
\]

The output tables contain:
\begin{itemize}
  \item the identities of the initial neighbors;
  \item the distances at each timepoint;
  \item summary values such as the mean and maximum distance across all tracked neighbors.
\end{itemize}

\section{Surface reconstruction}

\paragraph{Intuitive principle.}
The target surface is not a closed hull enclosing the full point cloud, but an open outer surface corresponding to the observed side of the embryo. The pipeline therefore prioritizes a microscope-axis-oriented reconstruction, analogous to a ``sheet'' stretched over the outermost points along the optical axis. This strategy limits artificial closure of concavities and the inclusion of internal points. If it fails, a 3D alpha-shape fallback method is used.

\subsection{Input data}

The surface is reconstructed independently at each timepoint from the cellular point cloud. The pipeline computes surfaces:
\begin{itemize}
  \item on raw coordinates;
  \item on stabilized coordinates;
  \item on interpolated timepoints for the viewer \emph{smooth} mode.
\end{itemize}

The default method is an open surface oriented along the microscope $z$ axis, modeled as a sheet of the form
\[
z=f(x,y).
\]

\subsection{Projection and estimation of local scale}

For each point
\[
\mathbf{x}_i=(x_i,y_i,z_i),
\]
we project onto the plane orthogonal to the microscope axis:
\[
\mathbf{u}_i=(x_i,y_i).
\]

From nearest-neighbor distances in this plane, a local spacing $s_i$ and a robust global spacing $s_0$ are estimated.

In the current implementation:
\begin{itemize}
  \item $s_i$ is the 75th percentile of projected distances to the $k=6$ nearest neighbors;
  \item $s_0$ is the 70th percentile of these distances at the global scale.
\end{itemize}

\subsection{Automatic choice of the outer side}

The code tests two hypotheses:
\begin{itemize}
  \item the outer surface lies on the side of maximal $z$ values;
  \item the outer surface lies on the side of minimal $z$ values.
\end{itemize}

For each hypothesis, a candidate sheet is extracted and a regularity score is computed:
\[
\mathrm{score}=
\mathrm{roughness}
+0.15\,\mathrm{spacing},
\]
with
\[
\mathrm{roughness}
=
\operatorname{median}\left(
\frac{|z_i-z_j|}{\|\mathbf{u}_i-\mathbf{u}_j\|_2}
\right),
\]
where $j$ runs over the projected neighbors of $i$, and
\[
\mathrm{spacing}
=
\operatorname{median}\big(\|\mathbf{u}_i-\mathbf{u}_j\|_2\big).
\]

The retained side is the one minimizing this score.

\subsection{Selection of outer points using a 2D grid}

We define a grid cell size
\[
h=1.05\,s_0.
\]

Four shifted grids are used:
\[
(0,0),\quad
\left(\frac{1}{2},0\right),\quad
\left(0,\frac{1}{2}\right),\quad
\left(\frac{1}{2},\frac{1}{2}\right).
\]

Within each grid cell, only the outermost point along the selected axis is retained:
\[
i^\star=\arg\max_{i\in \text{cell}} z_i
\quad\text{or}\quad
i^\star=\arg\min_{i\in \text{cell}} z_i.
\]

Vertices of the projected 2D convex hull are systematically added in order to preserve the lateral boundary.

\subsection{Triangulation}

On the selected points in the $(x,y)$ plane, a 2D Delaunay triangulation is constructed. Each projected triangle is then mapped to a 3D triangle on the sheet.

\subsection{Geometric filtering of triangles}

For a triangle $f=(i,j,k)$, we define:
\begin{align*}
e_{ij} &= \|\mathbf{u}_i-\mathbf{u}_j\|_2,\\
e_{ik} &= \|\mathbf{u}_i-\mathbf{u}_k\|_2,\\
e_{jk} &= \|\mathbf{u}_j-\mathbf{u}_k\|_2,\\
e_{\max}(f) &= \max(e_{ij},e_{ik},e_{jk}),\\
\bar{s}_f &= \frac{s_i+s_j+s_k}{3}.
\end{align*}

The 2D circumradius of the triangle is
\[
R_f=\frac{e_{ij}e_{ik}e_{jk}}{4A_f},
\]
where $A_f$ is the area of the projected triangle.

The axial variation of the triangle is
\[
\Delta z_f=
\max\big(
|z_i-z_j|,\,
|z_i-z_k|,\,
|z_j-z_k|
\big).
\]

A triangle is retained if the following three conditions are satisfied:
\begin{align*}
e_{\max}(f) &< 3.35\,\bar{s}_f,\\
R_f &< 2.95\,\bar{s}_f,\\
\Delta z_f &< 5.5\,\bar{s}_f.
\end{align*}

In addition, the code removes \emph{long and pointy} triangles. If $\theta_{\min}(f)$ is the smallest internal angle of the triangle, then the triangle is rejected if
\[
e_{\max}(f)>1.30\,\bar{s}_f
\]
and simultaneously
\[
\left(
\frac{e_{\max}(f)}{e_{\min}(f)}>3.4
\right)
\quad\text{or}\quad
\left(
\theta_{\min}(f)<22^\circ
\right).
\]

A local repair of small holes is then applied: a removed triangle may be reinserted if it remains close to the previous thresholds and already shares at least two edges with the retained mesh.

\subsection{Normal orientation}

The ordering of the vertices of each face is reversed when necessary so that the normal points globally toward the selected outer side.

\subsection{Fallback method: 3D alpha-shape}

If the oriented sheet cannot be constructed, the pipeline switches to a 3D \emph{alpha-shape}.

\paragraph{Preselection of surface points.}
Each point receives a score combining:
\begin{enumerate}[label=\arabic*.]
  \item a local directional asymmetry;
  \item a local planarity.
\end{enumerate}

For a point $i$, if $\mathbf{v}_{ij}$ are the unit vectors toward its $k$ neighbors:
\[
a_i=\frac{1}{k}\left\|
\sum_{j=1}^{k}\mathbf{v}_{ij}
\right\|_2.
\]

If $\sigma_{i,1}\ge \sigma_{i,2}\ge \sigma_{i,3}$ are the singular values of its centered local neighborhood:
\[
p_i=
1-\frac{\sigma_{i,3}}{\sigma_{i,1}+\varepsilon}.
\]

The combined score is
\[
s_i=0.55\,a_i+0.45\,p_i.
\]

Points whose $s_i$ belongs to the upper fraction of the cloud are retained as candidate surface points.

\paragraph{Alpha-shape.}
A 3D Delaunay tetrahedralization is computed. A tetrahedron is retained if its circumsphere radius is smaller than
\[
\alpha = 2.0\,s_0,
\]
where $s_0$ is a robust global spacing of the cloud.

Surface faces are the faces of tetrahedra appearing exactly once in the retained tetrahedron set. These faces are then filtered by long-edge and pointy-triangle criteria, analogously to the oriented sheet case.

\section{Surface density}

\paragraph{Intuitive principle.}
Density is defined as a local property of the surface, evaluated at its vertices. The idea is to measure how compact the surrounding cellular neighborhood is around each mesh vertex. The closer the nearest cells are to that vertex, the higher the density. The computation relies on a $k$ nearest-neighbor estimator, here with $k=6$, followed by a robust normalization for visualization.

\subsection{Adaptive mesh refinement}

Before density is computed, overly large triangles are subdivided in order to avoid artificial flat color patches.

For each triangle $f$, if
\[
e_{\max}(f) > 1.8\,\bar{s}_f,
\]
it is subdivided into four sub-triangles by inserting edge midpoints. This operation is repeated at most twice.

\subsection{Definition of density}

Density is computed at each vertex $\mathbf{v}$ of the final mesh. The $k=6$ nearest cells are searched in the full point cloud of the considered timepoint, and their distances are denoted
\[
d_1(\mathbf{v}),\dots,d_k(\mathbf{v}).
\]

We define the local mean distance
\[
\bar{d}(\mathbf{v})=\frac{1}{k}\sum_{m=1}^{k} d_m(\mathbf{v}),
\]
then the raw density
\[
\rho_{\mathrm{raw}}(\mathbf{v})=\frac{1}{\bar{d}(\mathbf{v})}.
\]

This definition yields a high density when cells are close to the considered vertex.

\subsection{Robust normalization}

Density is normalized within each timepoint using the 2\% and 98\% percentiles:
\[
p_2 = \operatorname{Perc}_{2\%}\big(\rho_{\mathrm{raw}}\big),
\qquad
p_{98} = \operatorname{Perc}_{98\%}\big(\rho_{\mathrm{raw}}\big).
\]

The normalized density is
\[
\rho(\mathbf{v})=
\operatorname{clip}\left(
\frac{\rho_{\mathrm{raw}}(\mathbf{v})-p_2}{p_{98}-p_2},
0,1
\right).
\]

In the exported file, this value is quantized as an unsigned 8-bit integer:
\[
\rho_{8\text{bits}}(\mathbf{v})=\lfloor 255\,\rho(\mathbf{v})\rfloor.
\]

This density is therefore a \emph{within-timepoint relative measure}. It is well suited to visualize spatial gradients, but should not be interpreted as an absolute density normalized across samples without additional recalibration.

\section{Data exported to the viewer}

\paragraph{Intuitive principle.}
The export is designed to clearly separate biological objects tracked over time from the geometric objects used for surface visualization. The viewer therefore receives both the cell trajectories, the lineage information, and the surface meshes enriched with their local density.

For each cell, the viewer receives:
\begin{itemize}
  \item raw coordinates;
  \item stabilized coordinates;
  \item existence timepoints;
  \item lineage information;
  \item neighborhood tables.
\end{itemize}

For each surface and each timepoint, the viewer receives:
\begin{itemize}
  \item the list of vertices;
  \item the list of faces;
  \item the per-vertex density;
  \item metadata describing the method used (\texttt{oriented\_sheet} or \texttt{alpha\_shape}) and the selected outer side.
\end{itemize}

\end{document}
